% ═══════════════════════════════════════════════════════════════════════════════
%  CAPÍTULO — Flyweight
% ═══════════════════════════════════════════════════════════════════════════════

\chapter{Flyweight}

\patroninfo{Estructural}{195}

% ─── Situación Problema ──────────────────────────────────────────────────────
\section{Situación Problema}

La plataforma de Airbnb almacena en memoria cientos de miles de objetos
\texttt{Propiedad} para servir búsquedas en tiempo real. Cada propiedad
puede declarar hasta cincuenta amenidades: Wi-Fi, piscina, cocina
equipada, aire acondicionado, estacionamiento, etc. Si cada objeto
\texttt{Propiedad} almacenara su propia instancia de cada amenidad —con
su nombre, icono, descripción y categoría— la memoria consumida se
multiplicaría de forma innecesaria, dado que la inmensa mayoría de
propiedades comparte exactamente las mismas amenidades con los mismos
datos descriptivos.

El problema se hace crítico cuando el motor de búsqueda carga en caché
los resultados para miles de usuarios concurrentes: el consumo de memoria
se dispara con objetos masivamente duplicados que son, en esencia,
idénticos.

% ─── Solución ────────────────────────────────────────────────────────────────
\section{Solución}

El patrón Flyweight separa el estado \emph{intrínseco} de una amenidad
—nombre, descripción, icono, categoría, que son inmutables y compartibles—
del estado \emph{extrínseco} —si una propiedad concreta la ofrece o no—.
Una \texttt{FabricaAmenidades} mantiene un mapa de objetos
\texttt{Amenidad} ya creados; cuando una propiedad declara que ofrece
Wi-Fi, recibe la misma instancia compartida que ya usan otras miles de
propiedades.

El resultado es una reducción drástica en el número de objetos en memoria:
en lugar de crear millones de instancias de \texttt{Amenidad} duplicadas,
el sistema reutiliza un conjunto acotado de objetos compartidos. Cada
propiedad solo guarda referencias a las amenidades relevantes, no copias
de sus datos.

% ─── Diagrama de clases ─────────────────────────────────────────────────────
\begin{figure}[H]
    \centering
    % \includegraphics[width=\textwidth]{imagenes/abstract_factory_diagrama.png}
    \caption{Diagrama de clases del patrón Abstract Factory}
\end{figure}


% ─── Roles de las Clases ────────────────────────────────────────────────────
\section{Roles de las Clases}

% Explicar la responsabilidad de cada clase/interfaz

\begin{itemize}
    \item \textbf{AbstractFactory:}
    \item \textbf{ConcreteFactory:}
    \item \textbf{AbstractProduct:}
    \item \textbf{ConcreteProduct:}
    \item \textbf{Client:}
\end{itemize}


% ─── Main ───────────────────────────────────────────────────────────────────
\section{Main}

% Explicar qué realiza el programa principal

\begin{lstlisting}[language=Java, caption=NombreClase]
 
\end{lstlisting}


% ─── Salida del Main ────────────────────────────────────────────────────────
\section{Salida del Main}

\begin{lstlisting}[language=Java, caption=NombreClase]
 
\end{lstlisting}


% ─── Clases ─────────────────────────────────────────────────────────────────
\section{Clases}

% ─── Clase 1 ────────────────────────────────────────────────────────────────
\subsection{NombreClase}

\begin{lstlisting}[language=Java, caption=NombreClase]
 
\end{lstlisting}


% ─── Clase 2 ────────────────────────────────────────────────────────────────
\subsection{NombreClase}

\begin{lstlisting}[language=Java, caption=NombreClase]
 
\end{lstlisting}


% ─── Clase 3 ────────────────────────────────────────────────────────────────
\subsection{NombreClase}

\begin{lstlisting}[language=Java, caption=NombreClase]
 
\end{lstlisting}


% ─── Conclusiones ───────────────────────────────────────────────────────────
\section{Conclusiones}

% Explicar:
% - Ventajas del patrón
% - Cuándo usarlo
% - Beneficios obtenidos
% - Posibles desventajas
% - Relación con principios SOLID